\documentclass[14pt, a4paper]{article}
\usepackage[T2A]{fontenc}
\usepackage[utf8]{inputenc}
\usepackage[russian]{babel}
\usepackage[left=3cm, right=1.5cm, top=2cm, bottom=2cm]{geometry}
\usepackage{indentfirst}
\usepackage{microtype}
\usepackage{array}
\usepackage{xcolor}
% \usepackage{times}
\usepackage{tempora}
\usepackage{ulem}
\usepackage{setspace}
\pagestyle{empty}

\makeatletter
\renewcommand\section{\@startsection {section}{1}{\parindent}%
  {-3.5ex \@plus -1ex \@minus -.2ex}%
  {0.5ex}%
  {\normalfont\Large\bfseries}}
\renewcommand\tableofcontents{%
  \begin{center}\Large\bfseries\contentsname\end{center}%
  \@starttoc{toc}%
}
\renewcommand*\l@section{\@dottedtocline{1}{0em}{1.5em}}
\makeatother

\begin{document}
\onehalfspacing

\begin{center}
    \textbf{ФЕДЕРАЛЬНОЕ ГОСУДАРСТВЕННОЕ АВТОНОМНОЕ ОБРАЗОВАТЕЛЬНОЕ} \\
    \textbf{УЧРЕЖДЕНИЕ ВЫСШЕГО ОБРАЗОВАНИЯ} \\
    \textbf{«МОСКОВСКИЙ АВИАЦИОННЫЙ ИНСТИТУТ} \\
    \textbf{(НАЦИОНАЛЬНЫЙ ИССЛЕДОВАТЕЛЬСКИЙ УНИВЕРСИТЕТ)»}
\end{center}

\begin{center}
    \bfseries Институт №8 «Компьютерные науки и \\
    \bfseries прикладная математика» \\
    \bfseries Кафедра 806 «Вычислительная математика \\
    \bfseries и программирование»
\end{center}

\vspace{3cm}

\begin{center}
    \Large\bfseries Отчет \\
    \vspace{0.5cm}
    \mdseries по лабораторной работе №3.1
\end{center}

\vfill

\noindent
\hfill
\begin{tabular}{ll}
Группа:       & М8О-307Б-21 \\
Студент:      & М. С. Ивченко \\
Преподаватель:& доц. В. В. Киндинова \\
Оценка:       & \underline{\makebox[3cm][c]{}} \\ [0.2cm]
Дата:         & \underline{\makebox[3cm][c]{}}
\end{tabular}

\vfill

\begin{center}
    Москва \\
    2026
\end{center}

\newpage

\setcounter{tocdepth}{2}
\tableofcontents
\newpage

\section*{Тема}
\addcontentsline{toc}{section}{Тема}
Цель данной лабораторной работы -- изучить восстановление функции по конечному набору ее значений. В работе строятся интерполяционные многочлены Лагранжа и Ньютона, а затем оценивается, насколько точно они приближают функцию в заданной точке.

\section*{Задание}
\addcontentsline{toc}{section}{Задание}
Используя таблицу значений $Y_i$ функции $y = f(x)$, вычисленных в точках $X_i$, $i = 0, \dots, 3$, построить интерполяционные многочлены Лагранжа и Ньютона, проходящие через точки $\{X_i, Y_i\}$. Вычислить значение погрешности интерполяции в точке $X^*$.

\section*{Теория}
\addcontentsline{toc}{section}{Теория}
% TODO: Описать интерполяционный многочлен Лагранжа, разделенные разности, многочлен Ньютона и оценку погрешности интерполяции.

\section*{Ход лабораторной работы}
\addcontentsline{toc}{section}{Ход лабораторной работы}
% TODO: Описать табличные данные варианта, построение многочленов Лагранжа и Ньютона, вычисление значения в точке X^* и погрешности.

\section*{Выводы}
\addcontentsline{toc}{section}{Выводы}
% TODO: Сформулировать выводы по лабораторной работе 3.1.

\end{document}
